% Part: lambda-calculus
% Chapter: lambda-definability
% Section: truth-values

\documentclass[../../../include/open-logic-section]{subfiles}

\begin{document}

\olfileid{lam}{ldf}{tvr}
\olsection{صداقتی قدراں تے تعلق}

خالص لامبڈا کیلکولس وچ صداقتی قدراں نوں ایویں لکھ سکدے آں۔
\begin{align*}
  \fn{true} & \ident \lambd[x][\lambd[y][x]]\\
  \fn{false} & \ident \lambd[x][\lambd[y][y]]
\end{align*}

صداقتی قدراں نوں \emph{چُنن والے فنکشوناں} دے روپ وچ لکھدے
آں، یعنی اوہ فنکشن جیہڑے دو دلیلاں لیندے نیں تے اوہناں وچوں
اک واپس دیندے نیں۔ صداقتی قدر $\fn{true}$ اپنی پہلی دلیل
چُندی اے تے $\fn{false}$ دوجی۔ مثلاً $\fn{true}\, M N$ دی
تقلیل ہمیشہ $M$ تک ہوندی اے، تے $\fn{false}\, M N$ دی
ہمیشہ~$N$ تک۔

\begin{defn}
اسیں تعلق $R \subseteq \Nat^n$ نوں !!{lambda definable}
آکھدے آں جے کوئی ٹرم~$R$ ایہو جیہا ہووے کہ
\begin{align*}
  R\, \num{n_1} \dots \num{n_k} & \bred \fn{true}
  \intertext{جدوں وی $R(n_1, \dots, n_k)$ قائم ہووے، تے}
  R\, \num{n_1} \dots \num{n_k} & \bred \fn{false}
\end{align*}
ہور صورت وچ۔
\end{defn}

مثلاً تعلق $\fn{IsZero} = \{0\}$ جیہڑا $0$ لئی تے صرف $0$
لئی قائم اے، ہِٹھلے ٹرم نال !!{lambda definable} اے۔
\[
  \fn{IsZero} \ident \lambd[n][n (\lambd[x][\fn{false}])\, \fn{true}].
\]
ایہ کیویں کم کردا اے؟ چرچ عدد دہراؤ والے فنکشن نیں جیہڑے
پہلی دلیل نوں $n$ واری دوجی اُتے لاؤندے نیں۔ سو پہلاں
قدر $\fn{true}$ رکھدے آں۔ تکرار دے ہر قدم اُتے پچھلے
قدم دا نتیجہ جو وی ہووے، $\fn{false}$ واپس دیندے آں۔
ایہ قدم مُڈھلی قدر اُتے $n$ واری لگے گا۔ نتیجہ $\fn{true}$
اُدوں تے صرف اُدوں ہووے گا جدوں ایہ قدم اک واری وی نہ لگے،
یعنی جدوں $n = 0$۔

صداقتی قدراں دی ایس نمائندگی توں اسیں ہور صداقتی فنکشن
وی متعین کر سکدے آں۔ نفی تے عطف لئی دو فنکشن ایہ نیں۔
\begin{align*}
  \fn{Not} & \ident \lambd[x][x\, \fn{false}\, \fn{true}]\\
  \fn{And} & \ident \lambd[x][\lambd[y][xy \,\fn{false}]]
\end{align*}
فنکشن «$\fn{Not}$» اک دلیل لیندا اے۔ ایہ $\fn{true}$ واپس
دیندا اے جے دلیل $\fn{false}$ ہووے، تے $\fn{false}$ واپس
دیندا اے جے دلیل~$\fn{true}$ ہووے۔ فنکشن «$\fn{And}$» دو
صداقتی قدراں دلیل وجوں لیندا اے، تے $\fn{true}$ اُدوں تے
صرف اُدوں واپس دینا چاہیدا اے جدوں دوویں دلیلاں~$\fn{true}$
ہون۔ جیویں اُتے دسیا اے، صداقتی قدراں چُنن والے فنکشن نیں۔
سو جدوں $x$ صداقتی قدر ہووے تے اینوں دو دلیلاں اُتے لائیے،
تاں پہلی دلیل نتیجہ ہووے گی جے $x$، $\fn{true}$ اے، تے ہور
صورت وچ دوجی دلیل۔ ہُن $\fn{And}$ اپنی دو دلیلاں $x$
تے $y$ لیندا اے تے $y$ تے $\fn{false}$ نوں اپنی پہلی
دلیل~$x$ اُتے لاؤندا اے۔ جے $x$ صداقتی قدر اے، تاں
نتیجہ~$y$ اے جے $x$، $\fn{true}$ اے، تے $\fn{false}$ اے
جے $x$، $\fn{false}$ اے۔ سانوں ایہو نتیجہ چاہیدا سی۔

دھیان دیو کہ ایتھے اسیں منّ کے چلدے آں کہ $\fn{And}$ دیاں
دلیلاں صرف صداقتی قدراں نیں۔ جے اینوں ہور ٹرم دِتّے جان،
تاں ہو سکدا اے کہ نتیجہ، یعنی نارمل روپ جے موجود ہووے، صداقتی
قدر نہ ہووے۔

\begin{prob}
  فنکشن $\fn{Or}$ تے $\fn{Xor}$ متعین کرو، جیہڑے شمولی تے
  اخراجی فصل دے صداقتی فنکشن ظاہر کردے نیں۔ ایس لئی صداقتی
  قدراں نوں $\lambd$-ٹرماں دے روپ وچ لکھن دا طریقہ ورتو۔
\end{prob}

\end{document}
